% Part II -- Python Foundations
% Chapter 5: Loops and Repetition

\h{Loops and Repetition}

Repeating similar instructions by hand creates long code and makes mistakes
more likely. A loop tells Python to repeat a block of instructions. The values
may come from a list, a dictionary, a numeric range, or a condition that
changes while the program runs.

\begin{NoteBox}
\textbf{Prerequisites.} You should understand lists, dictionaries, booleans,
comparisons, conditional statements, and indentation.
\end{NoteBox}

\begin{tcolorbox}[title=Learning outcomes,colframe=orange,colback=white,arc=2mm]
After completing this chapter, you should be able to:
\begin{itemize}
    \item repeat instructions with a \c{for} loop;
    \item use \c{range()} for a known number of repetitions;
    \item calculate a running total;
    \item iterate over dictionary keys and values;
    \item pair values with \c{enumerate()} and \c{zip()}; and
    \item control a basic \c{while} loop safely.
\end{itemize}
\end{tcolorbox}

\hh{Why loops are useful}

Suppose a list contains four room names. Repeating \c{print()} with four
indexes works, but the code must change whenever the list changes.

\code{c05_without_loop}{python}{
    rooms = ["Lobby", "Studio", "Office", "Cafe"]

    print(rooms[0])
    print(rooms[1])
    print(rooms[2])
    print(rooms[3])
}{Repeated instructions without a loop}

A loop expresses the real intention: display every item.

\code{c05_first_loop}{python}{
    rooms = ["Lobby", "Studio", "Office", "Cafe"]

    for room in rooms:
        print(room)
}{Display every item with a for loop}

Read the loop header as ``for each \c{room} in \c{rooms}.'' During each
repetition, \c{room} refers to the next item.

\hh{Understand the loop block}

The colon ends the loop header. Indented instructions form the loop body and
run once for each item. Instructions after the indentation ends run once after
the loop.

\code{c05_loop_block}{python}{
    rooms = ["Lobby", "Studio", "Office"]

    print("Room list")
    for room in rooms:
        print(room)
        print("Checked")

    print("All rooms checked")
}{Distinguish repeated and non-repeated instructions}

\begin{NoteBox}
\textbf{Stop and predict.} How many times will \c{Checked} appear? How many
times will \c{All rooms checked} appear?
\end{NoteBox}

\hh{Transform a value during each repetition}

The loop variable may be used in a calculation or f-string.

\code{c05_loop_calculation}{python}{
    areas_m2 = [24.0, 35.5, 18.0]

    for area_m2 in areas_m2:
        area_ft2 = area_m2 * 10.764
        print(f"{area_m2:.1f} m2 = {area_ft2:.1f} ft2")
}{Apply one calculation to every list item}

Each result is calculated from the current item. The original list is not
changed.

\hh{Repeat with range}

The \c{range()} function produces a sequence of integers. It is useful when
you know how many times to repeat.

\code{c05_range_basic}{python}{
    for number in range(5):
        print(number)
}{Count from zero through four}

The stop value is not included. \c{range(5)} produces 0, 1, 2, 3, and 4.

Provide a start and stop value when counting should begin elsewhere.

\code{c05_range_start_stop}{python}{
    for level in range(1, 5):
        print(f"Level {level}")
}{Count from one through four}

A third argument sets the step.

\code{c05_range_step}{python}{
    for grid_line in range(0, 25, 5):
        print(grid_line)
}{Count from zero by fives}

\hh{Calculate a running total}

An accumulator stores a result that changes during the loop. Begin at zero,
then add each value.

\code{c05_running_total}{python}{
    areas_m2 = [24.0, 35.5, 18.0]
    total_area_m2 = 0.0

    for area_m2 in areas_m2:
        total_area_m2 = total_area_m2 + area_m2

    print(f"Total area: {total_area_m2:.1f} m2")
}{Accumulate a total}

The shorter statement \c{total_area_m2 += area_m2} has the same purpose.

\code{c05_running_total_short}{python}{
    areas_m2 = [24.0, 35.5, 18.0]
    total_area_m2 = 0.0

    for area_m2 in areas_m2:
        total_area_m2 += area_m2

    average_area_m2 = total_area_m2 / len(areas_m2)
    print(f"Average area: {average_area_m2:.1f} m2")
}{Use an accumulator to calculate an average}

Keep the final calculation outside the loop because it is needed only once.

\hh{Combine loops and conditions}

An \c{if} statement inside a loop can select which items receive an action.

\code{c05_loop_with_if}{python}{
    areas_m2 = [24.0, 35.5, 18.0, 42.0]

    for area_m2 in areas_m2:
        if area_m2 >= 30.0:
            print(f"Large area: {area_m2:.1f} m2")
}{Display only values that meet a condition}

Indentation shows the structure: the \c{if} belongs to the loop, and its
\c{print()} belongs to the \c{if}.

\hh{Iterate over dictionaries}

Looping directly over a dictionary produces its keys.

\code{c05_dictionary_keys}{python}{
    room_areas = {
        "Lobby": 24.0,
        "Studio": 35.5,
        "Office": 18.0
    }

    for room_name in room_areas:
        print(room_name)
}{Iterate over dictionary keys}

Use \c{items()} when both the key and value are needed.

\code{c05_dictionary_items}{python}{
    room_areas = {
        "Lobby": 24.0,
        "Studio": 35.5,
        "Office": 18.0
    }

    for room_name, area_m2 in room_areas.items():
        print(f"{room_name}: {area_m2:.1f} m2")
}{Iterate over dictionary key--value pairs}

The two loop variables receive the key and value from each pair.

\hh{Number items with enumerate}

The \c{enumerate()} function provides an index and an item during each
repetition.

\code{c05_enumerate}{python}{
    rooms = ["Lobby", "Studio", "Office"]

    for number, room_name in enumerate(rooms, start=1):
        print(f"{number}. {room_name}")
}{Number list items starting at one}

Expected output:

\code{c05_enumerate_output}{text}{
    1. Lobby
    2. Studio
    3. Office
}{Numbered output from enumerate}

Use \c{enumerate()} instead of manually updating a counter when you need both
position and value.

\hh{Pair collections with zip}

The \c{zip()} function pairs corresponding items from two or more collections.

\code{c05_zip}{python}{
    room_names = ["Lobby", "Studio", "Office"]
    areas_m2 = [24.0, 35.5, 18.0]

    for room_name, area_m2 in zip(room_names, areas_m2):
        print(f"{room_name}: {area_m2:.1f} m2")
}{Pair room names and areas}

If the collections have different lengths, \c{zip()} stops at the end of the
shortest one. Check that related collections have matching lengths.

\hh{Repeat while a condition is true}

A \c{while} loop repeats as long as its condition remains true. Something in
the loop must eventually make the condition false.

\code{c05_first_while}{python}{
    count = 1

    while count <= 4:
        print(count)
        count += 1

    print("Finished")
}{A while loop with an updating counter}

Trace the counter carefully:
\begin{enumerate}
    \item \c{count} begins at 1.
    \item Python checks \c{count <= 4}.
    \item The body displays and increases \c{count}.
    \item When \c{count} becomes 5, the condition is false.
\end{enumerate}

Forgetting \c{count += 1} creates an infinite loop. Use
\texttt{Ctrl+C} in the terminal to stop a program that does not end.

\hh{Stop a loop with break}

The \c{break} statement exits the nearest loop immediately.

\code{c05_break}{python}{
    room_names = ["Lobby", "Studio", "Office", "Cafe"]
    target = "Office"

    for room_name in room_names:
        if room_name == target:
            print("Found the target")
            break

        print(f"Checked {room_name}")
}{Stop after finding a target}

Use \c{break} when continuing would serve no purpose.

\hh{Skip one repetition with continue}

The \c{continue} statement skips the rest of the current repetition and moves
to the next item.

\code{c05_continue}{python}{
    values = [24.0, None, 35.5, None, 18.0]

    for value in values:
        if value is None:
            continue

        print(f"Area: {value:.1f} m2")
}{Skip missing values}

Use \c{continue} only when it makes the loop easier to read than an additional
level of nesting.

\hh{Optional extension: a simple list comprehension}

After an ordinary loop is clear, a short list comprehension can express a
simple transformation.

\code{c05_ordinary_transform}{python}{
    areas_m2 = [20.0, 30.0, 40.0]
    doubled_areas = []

    for area_m2 in areas_m2:
        doubled_areas.append(area_m2 * 2)

    print(doubled_areas)
}{Build a transformed list with an ordinary loop}

The equivalent comprehension is:

\code{c05_list_comprehension}{python}{
    areas_m2 = [20.0, 30.0, 40.0]
    doubled_areas = [area_m2 * 2 for area_m2 in areas_m2]

    print(doubled_areas)
}{The same transformation as a list comprehension}

Prefer the ordinary loop when the transformation includes several steps or
would be difficult to explain in one short sentence.

\hh{Common mistakes}

\begin{itemize}
    \item \textbf{Forgetting the colon}: loop headers end with a colon.
    \item \textbf{Incorrect indentation}: only indented instructions repeat.
    \item \textbf{Using an index when the item is enough}: write
          \c{for room in rooms} instead of indexing unnecessarily.
    \item \textbf{Placing the final calculation inside the loop}: calculate once
          after accumulating all values.
    \item \textbf{Forgetting to update a while-loop variable}: the loop may never end.
    \item \textbf{Using collections of unequal length with zip}: extra values are ignored.
\end{itemize}

\hh{Practice ladder}

\hhh{Level 0 -- Read and predict}

\code{c05_practice_level0}{python}{
    names = ["A", "B", "C"]

    for name in names:
        print(name)

    print("Done")
}{Predict repeated and final output}

\hhh{Level 1 -- Modify}

Add two names to the list and change the message inside the loop.

\hhh{Level 2 -- Complete}

Complete the accumulator so the program displays the total.

\code{c05_practice_level2}{python}{
    areas_m2 = [20.0, 32.5, 18.0, 29.5]
    total_area_m2 = 0.0

    for area_m2 in areas_m2:
        total_area_m2 += area_m2

    print(f"Total: {total_area_m2:.1f} m2")
}{A completed accumulator}

\hhh{Level 3 -- Apply}

Given matching lists of four room names and four areas, use \c{zip()} to
display each pair. Display \c{Large} after a pair when its area is at least 30.

\textbf{Hints}
\begin{itemize}
    \item \textbf{Hint 1}: Use two loop variables with \c{zip()}.
    \item \textbf{Hint 2}: Put an \c{if} statement inside the loop.
    \item \textbf{Hint 3}: Indent the conditional message one level farther.
\end{itemize}

\hhh{Level 4 -- Challenge}

Extend Level 3 by counting how many areas are at least 30 and displaying the
count after the loop.

\hh{Practice solutions}

\textbf{Level 0:} The output is \c{A}, \c{B}, \c{C}, and then \c{Done}.

\textbf{Level 2:} The total is \c{100.0 m2}.

\textbf{Level 3}

\code{c05_solution_level3}{python}{
    room_names = ["Lobby", "Studio", "Office", "Cafe"]
    areas_m2 = [24.0, 35.5, 18.0, 31.0]

    for room_name, area_m2 in zip(room_names, areas_m2):
        print(f"{room_name}: {area_m2:.1f} m2")
        if area_m2 >= 30.0:
            print("Large")
}{One solution to Level 3}

\textbf{Level 4}

\code{c05_solution_level4}{python}{
    room_names = ["Lobby", "Studio", "Office", "Cafe"]
    areas_m2 = [24.0, 35.5, 18.0, 31.0]
    large_count = 0

    for room_name, area_m2 in zip(room_names, areas_m2):
        print(f"{room_name}: {area_m2:.1f} m2")
        if area_m2 >= 30.0:
            print("Large")
            large_count += 1

    print(f"Large room count: {large_count}")
}{One solution to Level 4}

\hh{Part II checkpoint}

\begin{enumerate}
    \item What does the loop variable represent during one repetition?
    \item Which instructions repeat in a \c{for} loop?
    \item Does \c{range(5)} include 5?
    \item Where should an accumulator be initialized?
    \item When should a final average be calculated?
    \item What does \c{dictionary.items()} provide?
    \item When is \c{enumerate()} useful?
    \item What happens when collections passed to \c{zip()} have different lengths?
    \item What must happen inside a \c{while} loop to ensure it ends?
    \item What is the difference between \c{break} and \c{continue}?
\end{enumerate}

\hh{Checkpoint answers}

\begin{enumerate}
    \item The current item from the collection.
    \item The indented instructions in the loop body.
    \item No. It produces values from 0 through 4.
    \item Before the loop.
    \item After the loop has processed all values.
    \item Key--value pairs.
    \item When both an item's position and value are needed.
    \item It stops when the shortest collection ends.
    \item A value or condition must change so the condition eventually becomes false.
    \item \c{break} exits the loop; \c{continue} skips to the next repetition.
\end{enumerate}

\begin{tcolorbox}[title=Part II summary,colframe=orange,colback=white,arc=2mm]
You can now store values in variables and collections, calculate results, make
decisions, and repeat instructions. These ideas form the foundation for writing
larger and more reliable Python programs in the following chapters.
\end{tcolorbox}

